\documentclass[11pt]{article}
\usepackage[a4paper, margin=2cm]{geometry}
\usepackage[utf8]{inputenc}
\usepackage[spanish]{babel}
\usepackage{amsmath, amssymb}
\usepackage{tikz}
\usepackage{pgfplots}
\pgfplotsset{compat=1.18}
\usepackage{enumitem}

\pagestyle{empty}

% ---------- FIGURAS (dato extra, ayudan a visualizar) ----------
\newcommand{\poligonoInscripto}[2]{%
% #1 = numero de lados, #2 = etiqueta
\begin{tikzpicture}[scale=1]
  \def\R{1.4}
  \draw[thick] (0,0) circle (\R);
  \draw[thick] (90:\R) \foreach \i in {1,...,#1} { -- (90+\i*360/#1:\R) };
  \node at (0,-1.9) {\footnotesize #2};
\end{tikzpicture}%
}

\begin{document}

% \shorthandoff{<>} es OBLIGATORIO acá (justo después de \begin{document},
% no sirve en el preambulo): babel-spanish activa "<" y ">" como caracteres
% activos para las comillas angulares («»), y \selectlanguage (que corre
% automáticamente al entrar al documento) los reactiva pisando cualquier
% \shorthandoff puesto antes en el preámbulo. Con "<" y ">" activos,
% cualquier \draw[->], \draw[<-] o \draw[<->] de TikZ rompe la compilación
% con el error "Argument of \language@active@arg> has an extra }."
% No sacar esta línea.
\shorthandoff{<>}

% La minipage tiene ALTURA FIJA (\textheight): esto es lo que hace que la
% práctica ocupe la hoja completa sin importar cuántos ejercicios tenga.
% Los "plus 1fill" en itemsep y en el \vspace antes de "Figuras" son glue
% elástico: LaTeX reparte el espacio sobrante entre esos puntos en el
% momento de compilar. Esto es INTENCIONAL y OBLIGATORIO: no reemplaces el
% itemsep ni el vspace por valores fijos sin "plus", y no le saques la
% altura fija a la minipage.
\begin{minipage}[t][\textheight]{\linewidth}
\begin{center}
{\Large \textbf{Geometría: Polígonos regulares inscriptos}}\\[2pt]
\rule{7cm}{0.5pt}
\end{center}

\vspace{0.5cm}

\begin{enumerate}[leftmargin=1.2cm, itemsep=0.5cm plus 1fill]

\item Un triángulo equilátero está inscripto en una circunferencia de radio $6\,\text{cm}$. Calculá la longitud del lado del triángulo.

\item Un hexágono regular está inscripto en una circunferencia de radio $8\,\text{cm}$. Calculá el perímetro del hexágono.

\item Un pentágono regular está inscripto en una circunferencia de radio $5\,\text{cm}$. Calculá la longitud aproximada de uno de sus lados.

\item Calculá la apotema de un hexágono regular inscripto en una circunferencia de radio $10\,\text{cm}$.

\item Un triángulo equilátero inscripto en una circunferencia tiene lado $L = 9\,\text{cm}$. Calculá el radio de la circunferencia circunscripta.

\item Calculá el área de un hexágono regular inscripto en una circunferencia de radio $6\,\text{cm}$.

\item Un pentágono regular tiene lado $L = 7\,\text{cm}$. Calculá el radio de la circunferencia en la que está inscripto.

\item Calculá el área de un pentágono regular inscripto en una circunferencia de radio $9\,\text{cm}$.

\item Un octógono regular está inscripto en una circunferencia de radio $4\,\text{cm}$. Calculá la longitud de su lado.

\item Comparando un triángulo, un pentágono y un hexágono inscriptos en una misma circunferencia de radio $r$, ¿cuál de los tres tiene mayor perímetro? Justificá sin hacer los cálculos numéricos.

\end{enumerate}

\vspace{0.8cm plus 1fill}

\begin{center}
{\large \textbf{Figuras}}\\[0.4cm]
\begin{tikzpicture}
  \node at (0,0)  {\poligonoInscripto{3}{Triángulo}};
  \node at (4,0)  {\poligonoInscripto{5}{Pentágono}};
  \node at (8,0)  {\poligonoInscripto{6}{Hexágono}};
\end{tikzpicture}
\end{center}
\end{minipage}

\newpage
\begin{center}
{\Large \textbf{Respuestas}}\\[2pt]
\rule{6cm}{0.5pt}
\end{center}

\begin{enumerate}[leftmargin=1.2cm, itemsep=0.3cm]
\item $L = 6\sqrt{3} \approx 10{,}39\,\text{cm}$
\item Perímetro $= 48\,\text{cm}$ (lado $= 8\,\text{cm}$)
\item $L \approx 5{,}88\,\text{cm}$
\item Apotema $\approx 8{,}66\,\text{cm}$
\item $R = 3\sqrt{3} \approx 5{,}20\,\text{cm}$
\item Área $\approx 93{,}53\,\text{cm}^2$
\item $R \approx 5{,}96\,\text{cm}$
\item Área $\approx 173{,}21\,\text{cm}^2$
\item $L \approx 3{,}06\,\text{cm}$
\item El triángulo, ya que a mayor cantidad de lados el polígono se aproxima más a la circunferencia y cada lado es más corto; con menos lados, cada lado subtiende un arco mayor.
\end{enumerate}

\end{document}